\chapter{Software de supervisión y análisis}
\label{chap:software}

\section{Visión general}

El software del ordenador consta de un servidor en Python (\fichero{web\_monitor.py}), dos páginas
web (\fichero{static/index.html} para la operación en vivo y \fichero{static/ensayos.html} para el
análisis) y varias herramientas auxiliares. El servidor se ejecuta con

\begin{lstlisting}[language=bash,numbers=none]
python web_monitor.py [puerto_serie]      # http://localhost:8080
\end{lstlisting}

y detecta solo el puerto del microcontrolador si no se indica.

\section{Servidor}

\subsection{Estructura concurrente}

El servidor combina un hilo del sistema operativo para el puerto serie con el bucle asíncrono de
FastAPI (figura~\ref{fig:servidor}).

\begin{figure}[H]
\centering
\begin{tikzpicture}[font=\small,>=Stealth,
  box/.style={draw,rounded corners=2pt,align=center,minimum height=10mm,inner sep=4pt,fill=white}]
  \node[box,minimum width=28mm] (usb) {Puerto USB\\ \scriptsize pySerial};
  \node[box,minimum width=34mm,right=10mm of usb,fill=blue!6] (hilo) {Hilo serie\\ \scriptsize decodifica · integra · graba};
  \node[box,minimum width=24mm,right=10mm of hilo] (cola) {deque\\ \scriptsize maxlen = 100};
  \node[box,minimum width=30mm,right=10mm of cola,fill=green!6] (bc) {Difusor asíncrono\\ \scriptsize FastAPI / Uvicorn};
  \node[box,minimum width=24mm,below=10mm of bc] (ws) {Navegadores\\ \scriptsize WebSocket};
  \node[box,minimum width=24mm,below=10mm of hilo] (csv) {cycles/*.csv};
  \node[box,minimum width=30mm,below=10mm of cola] (cmd) {\_handle\_command\\ \scriptsize valida · empaqueta};
  \draw[->,thick] (usb) -- (hilo);
  \draw[->,thick] (hilo) -- (cola);
  \draw[->,thick] (cola) -- (bc);
  \draw[->,thick] (bc) -- (ws);
  \draw[->,thick] (hilo) -- (csv);
  \draw[->,thick] (ws) -- (cmd);
  \draw[->,thick] (cmd.south) -- ++(0,-5mm) -| node[pos=0.25,below,font=\scriptsize]{tramas binarias} (usb.south);
\end{tikzpicture}
\caption{Estructura del servidor de supervisión.}
\label{fig:servidor}
\end{figure}

\paragraph{Hilo serie.} Abre el puerto (probando \codigo{/dev/ttyACM*}, \codigo{/dev/ttyUSB*},
\codigo{/dev/cu.usbmodem*} y \codigo{COM*}), vacía el búfer y entra en un bucle que lee tramas,
comprueba el CRC y las decodifica con \codigo{struct}. Cualquier excepción cierra el puerto y
provoca un reintento al cabo de un segundo, lo que cubre la desconexión física y el reinicio del
microcontrolador al reprogramarlo (RNF4). Por cada muestra de telemetría:
\begin{itemize}
  \item mide la frecuencia de muestreo contando las muestras del último segundo \emph{con el reloj
        del dispositivo}, no con el del ordenador;
  \item detecta que el reloj del microcontrolador ha vuelto a cero (reinicio) y en ese caso reinicia
        la ventana de medida y la integración. Sin esta comprobación, la ventana no se podaba nunca y
        la interfaz llegó a mostrar ``\SI{3813}{\hertz}'', que en realidad era el número de muestras
        acumuladas;
  \item integra la posición a partir de la velocidad, descartando saltos de tiempo anómalos;
  \item si hay un ensayo en grabación, escribe la fila correspondiente;
  \item deposita un mensaje JSON en la cola.
\end{itemize}

\paragraph{Cola y difusor.} La comunicación entre el hilo y el bucle asíncrono usa una
\codigo{deque} de la biblioteca estándar con un único productor y un único consumidor, segura sin cerrojos en
CPython para \codigo{append} y \codigo{popleft}. La longitud máxima de 100 mensajes acota la memoria
si no hay ningún navegador conectado: los mensajes más antiguos se descartan. El difusor envía cada
mensaje a todos los clientes, elimina los que han cerrado la conexión y, cuando la cola está vacía,
duerme \SI{20}{\milli\second}.

\paragraph{Estado compartido.} Los valores que necesitan varios contextos (posición, calibración,
estado del \emph{homing}) se guardan en un diccionario protegido por un \codigo{threading.Lock}.

\paragraph{Registro.} Todas las tramas, órdenes y confirmaciones se registran con marca de tiempo en
consola y en \fichero{web\_monitor.log}. El registro ha sido la principal herramienta de diagnóstico
de los problemas de comunicación.

\subsection{Grabación de ensayos}

La grabación es automática y la dirigen los paquetes de ciclo: cuando la fase pasa a ``hacia B'' o
``volviendo a A'' y no hay fichero abierto, se crea \fichero{cycles/AAAAMMDD-HHMMSS\_ciclo.csv}; cuando
la fase pasa a completado, abortado o reposo, se cierra. Cada fichero tiene una cabecera con las
constantes del ensayo (calibración de la célula, fuerza en el punto de tara, repeticiones y punto B)
y una fila por muestra (formato completo en el anexo~\ref{anx:csv}).

Se graban deliberadamente las \emph{cuentas brutas} de la célula y no newtons. La calibración puede
cambiar después del ensayo (de hecho, la de este trabajo es provisional), y con los datos brutos
basta con aplicar la nueva constante (RNF5).

\section{Página de supervisión}

La página principal (\fichero{index.html}) se organiza en tres zonas:

\begin{enumerate}
  \item \textbf{Estado}: indicador de conexión, estado de la máquina y código de error, modo, estado del
        servo y de la célula, tiempo del dispositivo y frecuencia de muestreo, finales de carrera con
        fase del \emph{homing} y recorrido medido, y una tarjeta destacada con la fuerza sobre la
        carga.
  \item \textbf{Controles en pestañas}:
  \begin{itemize}
    \item \emph{Servo}: dirección Modbus, inicialización, movimiento manual, \emph{homing}, parada,
          puesta a cero de la posición y apagado; velocidad, desplazamiento del origen y velocidad de
          \emph{homing}.
    \item \emph{Célula 1}: calibración (cuentas/N), fuerza en el punto de tara y tara. En la
          configuración avanzada se accede a la palabra B\_Config (ganancia de ×1,5 a ×192, cero del
          A/D, polaridad, puente completo o medio, integración larga, alimentación del puente y
          anulación de desequilibrio) y a ZMDI\_Config1 (reloj de 4 o \SI{1}{\mega\hertz}, tasa de
          conversión y modo continuo o a demanda), con el botón de programar. La configuración
          grabada se consulta automáticamente la primera vez que se abre la pestaña con el eje
          parado.
    \item \emph{Ciclo A$\leftrightarrow$B}: punto B (final de carrera o posición), velocidad y
          repeticiones; arranque y parada; lista de ensayos y enlace al visor.
    \item \emph{Cálculo teórico}: esquema del mecanismo y curvas de tensión, par y ángulo para un
          tambor de radio constante, con comparación con el par medido.
  \end{itemize}
  \item \textbf{Gráficas}: velocidad consigna y real con los finales de carrera y marcas verticales en
        los instantes en que se enviaron órdenes, y fuerza frente al tiempo.
\end{enumerate}

\section{Visor de ensayos}

El visor (\fichero{ensayos.html}) lista los ensayos grabados y, al abrir uno, muestra:

\begin{itemize}
  \item un \textbf{resumen} con número de muestras, duración, cadencia, repeticiones y rangos de
        posición, par y fuerza, y una tabla por repetición;
  \item \textbf{fuerza frente a posición}, con una traza por repetición, donde se ven directamente la
        ida, la vuelta y la histéresis;
  \item \textbf{fuerza y posición frente al tiempo}, con la posibilidad de superponer la curva
        teórica del capítulo~\ref{chap:modelo} para un tambor cilíndrico o una caracola lineal. Se
        introducen las cotas, la masa, el ángulo de la barra en el origen, la reducción y el sentido
        de giro;
  \item \textbf{par frente a posición y frente al tiempo}, con el par teórico $\tau = T\,r(\varphi)$;
  \item \textbf{velocidad consigna y real} con los finales de carrera.
\end{itemize}

Al superponer la curva teórica, el visor comprueba dos consistencias que han resultado muy útiles
para detectar datos de entrada erróneos:
\begin{enumerate}
  \item \emph{Escala}: compara el cable que paga el tambor durante el ensayo con el que puede absorber
        el mecanismo desde el ángulo inicial hasta la vertical en ese sentido. Si difieren más de un
        \SI{15}{\percent}, avisa y sugiere la reducción o los radios que lo harían cuadrar.
  \item \emph{Ángulo}: con la ecuación~\eqref{eq:angulo} calcula el ángulo de la barra en el origen,
        en el extremo y al final del ensayo, y avisa si el modelo exigiría superar la vertical o si el
        ángulo inicial introducido no es alcanzable.
\end{enumerate}

Por último, ofrece un \textbf{ajuste por mínimos cuadrados} de la recta que convierte las cuentas de
la célula en la fuerza teórica, $F = a\cdot\text{cuentas} + b$, con
\begin{equation}
  a = \frac{n\sum x_iy_i - \sum x_i\sum y_i}{n\sum x_i^2 - (\sum x_i)^2}, \qquad
  b = \frac{\sum y_i - a\sum x_i}{n},
\end{equation}
e informa del coeficiente de determinación $R^2$ y de la \emph{cobertura}, el porcentaje de muestras
donde la geometría tiene solución. Un $R^2$ alto con poca cobertura no garantiza un ajuste válido, y
por eso la interfaz lo señala.

\section{Herramientas auxiliares}

\begin{itemize}
  \item \fichero{leer\_ciclo.py}: resumen de un ensayo en consola, usando sólo la biblioteca estándar,
        y exportación sin cabecera para hojas de cálculo.
  \item \fichero{monitor.py} y \fichero{monitor\_posicion.py}: monitores de terminal de las primeras
        versiones del protocolo, con control por teclado.
  \item \fichero{test\_servo\_modbus.py}: barrido de velocidades y direcciones Modbus con un adaptador
        USB-RS485 y un CRC-16 propio, útil para identificar la configuración del accionamiento sin el
        microcontrolador.
  \item \fichero{memoria/figuras/generar\_figuras.py}: genera todas las figuras de resultados y las
        cifras de esta memoria a partir de los ficheros de ensayo y del modelo.
\end{itemize}

\section{Observaciones de calidad}

Durante la redacción de esta memoria se revisó el software y se identificaron los siguientes puntos
de mejora, que no afectan a los resultados:
\begin{itemize}
  \item \fichero{requirements.txt} sólo declara \codigo{pyserial}; faltan \codigo{fastapi} y
        \codigo{uvicorn[standard]}, que se instalan a mano.
  \item La cabecera del CSV decide si B es ``posición'' o ``final de carrera'' según el objetivo en
        revoluciones sea mayor que cero, no según el modo realmente pedido. Como la interfaz envía
        siempre el objetivo (1 rev por defecto), todos los ensayos de este trabajo aparecen como
        ``posición'' aunque se hicieron hasta el final de carrera. Los datos no se ven afectados; la
        corrección es incluir el modo en el paquete de ciclo.
  \item El servidor escucha en todas las interfaces (\codigo{0.0.0.0}) sin autenticación, así que
        cualquier equipo de la red puede mover la máquina. En un entorno compartido debería
        escuchar sólo en \codigo{127.0.0.1} o exigir credenciales.
  \item La posición se integra dos veces, en el firmware y en el servidor, con el mismo método. El
        CSV guarda la del servidor. Sería más limpio transmitir la del firmware en cada muestra.
  \item Algunos comentarios han quedado desfasados respecto al código: la cabecera del firmware
        todavía describe la versión de doble núcleo y el tamaño de muestra de 12 bytes.
\end{itemize}
